% Created (UTC): 2026-06-25T19:01:26Z
% Repository HEAD: 469ff48bb80ab3c4e221417ab9ac2934d8f3d22a
\documentclass[11pt]{article}
\usepackage[margin=1in]{geometry}
\usepackage{amsmath,amssymb,amsthm,mathtools}
\usepackage{microtype}
\emergencystretch=3em \hbadness=4000
\usepackage[colorlinks=true,linkcolor=blue,citecolor=blue,urlcolor=blue]{hyperref}
\DeclareMathOperator{\Li}{Li}
\newcommand{\Cl}{\operatorname{Cl}_2}
\theoremstyle{plain}\newtheorem{proposition}{Proposition}
\theoremstyle{remark}\newtheorem{remark}[proposition]{Remark}
\title{Where the Herglotz rational values meet the golden ladders\\ and the Clausen lattice}
\author{Vladimir Reshetnikov}
\date{2026-06-25}
\begin{document}
\maketitle
\begin{center}\small
Research note. Created (UTC) 2026-06-25T19:01:26Z; repository \texttt{HEAD}
\texttt{469ff48bb80ab3c4e221417ab9ac2934d8f3d22a}. Companion to
\texttt{herglotz-rational-values}; bridges to \texttt{golden-polylog-ladders} and
\texttt{clausen-roots-of-unity}. Reproducer
\texttt{src/PolyLog/tools/scratch/herglotz/bridges.py}.
\end{center}

\begin{abstract}
The general formula $F(2/q)-F(1)=\dots-\tfrac12\sum\Li_2(-\cot^2\tfrac{\pi k}q)$ introduces dilogarithms at
algebraic arguments; for $q=5$ these live in $\mathbb Q(\sqrt5)$, the field of our golden polylogarithm
ladders, and the Herglotz \emph{derivative} law introduces a Clausen sum $W(q)$ that lives in our
conductor-$q$ Clausen lattice. We make both connections explicit. The golden-argument dilogarithms of
$F(2/5)$ split into a symmetric part that reduces to $\Li_2(\tfrac15)$ plus golden constants and an
antisymmetric part that is a genuinely new $\mathbb Q(\sqrt5)$ real-dilogarithm constant; and the derivative
Clausen sum $W(q)$ is an element of the conductor-$q$ Clausen lattice over the real cyclotomic field
$\mathbb Q(\zeta_q)^+$, with clean closed forms whenever $\cot\frac{2\pi k}q\in\mathbb Q(\zeta_q)^+$. All
identities are verified to $\sim10^{-80}$.
\end{abstract}

\section{The golden bridge ($q=5$)}
For $q=5$ the two algebraic arguments are $-\cot^2\tfrac{\pi}{5}=-\tfrac{5+2\sqrt5}{5}$ and
$-\cot^2\tfrac{2\pi}{5}=-\tfrac{5-2\sqrt5}{5}$, both in $\mathbb Q(\sqrt5)$. Comparing our
$F(2/5)$ formula with Radchenko--Zagier's eq.~(24) (which uses the rational $\Li_2(4/5)$) gives a
dilogarithm functional equation tying the two:
\begin{proposition}\label{prop:bridge}
With $\ell_k:=\log^2(2\sin\tfrac{\pi k}5)$,
\[
  \Li_2\!\Bigl(\tfrac45\Bigr)+\Li_2\!\Bigl(-\cot^2\tfrac{\pi}{5}\Bigr)+\Li_2\!\Bigl(-\cot^2\tfrac{2\pi}{5}\Bigr)
  =2\log^2 2-2\log2\log\tfrac25-6(\ell_1+\ell_2).
\]
\end{proposition}
Eliminating the rational dilogarithm gives a closed form for the symmetric golden-argument sum:
\begin{proposition}\label{prop:goldensum}
With $\varphi=\tfrac{1+\sqrt5}2$,
\[
  \Li_2\!\Bigl(-\cot^2\tfrac{\pi}{5}\Bigr)+\Li_2\!\Bigl(-\cot^2\tfrac{2\pi}{5}\Bigr)
  =\Li_2\!\Bigl(\tfrac15\Bigr)-\frac{\pi^2}{6}-3\log^2\varphi+\tfrac14\log^2 5 .
\]
\end{proposition}
Substituting Proposition~\ref{prop:goldensum} into the $F(2/5)$ theorem re-expresses that Herglotz value
entirely through $\Li_2(\tfrac15)$ and golden constants --- the same constants the golden ladders reach for
from the other side.

\begin{remark}[A new $\mathbb Q(\sqrt5)$ constant]
Only the \emph{symmetric} combination reduces. The antisymmetric part
$X:=\tfrac12\bigl(\Li_2(-\cot^2\tfrac{\pi}{5})-\Li_2(-\cot^2\tfrac{2\pi}{5})\bigr)$ admits no integer relation
over $\{\pi^2,\log^2\varphi,\log^2 5,\log^2 2,\Li_2(\tfrac15),\Li_2(\tfrac25),\text{log products}\}$ (PSLQ
with a Champernowne canary, $100$ digits, $\mathrm{maxcoeff}=10^5$): it is a genuinely new real-dilogarithm
constant attached to $\mathbb Q(\sqrt5)$. (Consistently, $X$ is not balanced as a Bloch-group element, so
Borel's vanishing of $K_3(\mathbb Q(\sqrt5))\otimes\mathbb Q$ does not force it to reduce.)
\end{remark}

\section{The Clausen bridge (the derivative law)}
The derivative law $\tfrac2q F'(2/q)-(1+\log2)=\tfrac{q^2-4}{24q}\pi^2+W(q)$ carries the Clausen sum
$W(q)=\sum_{k=1}^{q-1}\cot\tfrac{2\pi k}q\,\Cl(\tfrac{2\pi k}q)$. Unlike the \emph{value} $F(2/q)$, this does
not reduce to logarithms: $W(q)$ is a $\mathbb Q(\zeta_q)^+$-linear combination of the conductor-$q$ Clausen
values $\Cl(\tfrac{2\pi k}q)$ --- i.e.\ an element of our conductor-$q$ Clausen lattice over the real
cyclotomic field. The coefficients $\cot\tfrac{2\pi k}q$ are rational over $\mathbb Q(\zeta_q)^+$ exactly when
$\sin\tfrac{2\pi k}q\in\mathbb Q(\zeta_q)^+$; in those cases $W(q)$ has a clean closed form, e.g.
\[
  W(8)=2\,\Cl\!\Bigl(\tfrac{\pi}4\Bigr)-2\,\Cl\!\Bigl(\tfrac{3\pi}4\Bigr),\qquad
  W(12)=\frac{22\sqrt3}{9}\Bigl(\Cl\!\Bigl(\tfrac{\pi}6\Bigr)-\Cl\!\Bigl(\tfrac{5\pi}6\Bigr)\Bigr).
\]
For $q=5,7,9,11$ the cotangents lie in a strictly larger field (e.g.\ $\cot\tfrac{2\pi}5\notin\mathbb Q(\sqrt5)$),
so no rational- or $\mathbb Q(\sqrt5)$-coefficient reduction exists, while membership in the conductor-$q$
Clausen lattice over $\mathbb Q(\zeta_q)^+$ persists. So differentiating $F$ at a rational moves from the
real-dilogarithm world of $F(2/q)$ into the imaginary-dilogarithm (Clausen) lattice we mapped at roots of
unity.

\section*{Verification}
Propositions~\ref{prop:bridge}--\ref{prop:goldensum}, the irreducibility of $X$, and the $W(8),W(12)$
closed forms are checked to $\sim10^{-80}$ in
\texttt{src/PolyLog/tools/scratch/herglotz/bridges.py}.
\end{document}
